\documentclass{article}

\title{Measuring Runtime, CPU, RAM, and Other Resource Usage in Python}
\date{}
\begin{document}
    \maketitle
    
    
    \section{Goal}
        Measure execution time and resource usage (CPU, RAM, and optionally I/O) from the start of a program until the end, and optionally retrieve the measurements at any intermediate point. Persist the collected data to disk (e.g., at program exit or at checkpoints).
    
    
    \section{Time Measurement}
        
        \subsection{Wall-Clock Time (Elapsed Real Time)}
            Use \texttt{time.perf\_counter()} to measure elapsed real time.
            It is suitable for benchmarking and includes waiting time (e.g., I/O waits and sleep).
        
        \subsection{CPU Time (Process CPU Consumption)}
            Use \texttt{time.process\_time()} to measure CPU time consumed by the current process.
            It excludes time when the process is not executing on CPU.
    
    
    \section{Memory Measurement (RAM)}
        
        \subsection{Python-Level Allocations (Inside the Python Interpreter)}
            Use \texttt{tracemalloc} (standard library) to track memory allocations attributed to Python code.
            This is useful to understand where Python memory grows over time and to take snapshots at checkpoints.
            It does not directly represent the full resident set size (RSS) of the process.
        
        \subsection{Process-Level Memory on Unix Systems}
            On Unix-like systems, the \texttt{resource} module can provide process resource statistics.
            A common metric is \texttt{ru\_maxrss}, which reports the maximum resident set size observed.
            Units can be platform-dependent (on Linux it is commonly KiB).
        
        \subsection{Comprehensive Process Metrics}
            The \texttt{psutil} library is commonly used for accurate and convenient process metrics.
            It can report RSS/VMS memory, CPU usage, thread counts, and I/O counters for the current process (via its PID).
            For most practical needs involving CPU and RAM over time, \texttt{psutil} is the most straightforward option.
    
    
    \section{Collecting Data Over the Entire Runtime}
        
        \subsection{Checkpoint-Based Collection}
            Record a measurement at specific points (manual checkpoints).
            This is appropriate when only a few milestones matter (e.g., after loading data, after training, after evaluation).
        
        \subsection{Periodic Sampling}
            Run a lightweight sampler that records metrics every N seconds throughout execution.
            This enables computing overall summaries such as peak memory, average CPU usage, and time-series plots of resource usage.
    
    
    \section{Persisting Measurements to Disk}
        
        \subsection{Common Output Formats}
            CSV is convenient for quick inspection and plotting.
            JSON Lines (JSONL) is robust for long-running logs, with one JSON record per line.
            SQLite is useful when you want structured storage and queries over measurements.
        
        \subsection{End-of-Program Summary}
            In addition to raw samples, write a summary at program termination (e.g., total elapsed time, total CPU time, peak RSS, and any relevant aggregates derived from samples).
    
    
    \section{Practical Recommendation}
        For measuring CPU and RAM from start to finish and retrieving values at any moment, combine:
        Wall-clock timing via \texttt{time.perf\_counter()} with process metrics via \texttt{psutil}.
        For diagnosing Python allocation growth and identifying sources of memory increase within Python code, add \texttt{tracemalloc} snapshots at selected checkpoints.

\end{document}
